\documentclass[11pt,a4paper]{article}
% ==============================================================
%  Préambule commun aux fiches -- Advanced Probabilities (M1)
%  Usage, dans chaque fiche :
%     \documentclass[11pt,a4paper]{article}
%     % ==============================================================
%  Préambule commun aux fiches -- Advanced Probabilities (M1)
%  Usage, dans chaque fiche :
%     \documentclass[11pt,a4paper]{article}
%     % ==============================================================
%  Préambule commun aux fiches -- Advanced Probabilities (M1)
%  Usage, dans chaque fiche :
%     \documentclass[11pt,a4paper]{article}
%     \input{../preambule.tex}
%     \fiche{C5}{Vecteurs gaussiens}{Gaussian vectors}
% ==============================================================
\usepackage[T1]{fontenc}
\usepackage[utf8]{inputenc}
\usepackage[french]{babel}
\usepackage{lmodern}
\usepackage[a4paper,margin=1.9cm,top=1.7cm,bottom=1.9cm]{geometry}
\usepackage{amsmath,amssymb,amsthm,mathtools}
\usepackage[inline]{enumitem}
\usepackage{booktabs,array,tabularx}
\usepackage{xcolor}
\usepackage[most]{tcolorbox}
\usepackage{fancyhdr}
\usepackage{hyperref}
\hypersetup{colorlinks=true,linkcolor=blue!50!black,urlcolor=blue!50!black}

\setlist{itemsep=1pt,topsep=3pt,parsep=0pt}
\setlength{\parindent}{0pt}
\setlength{\parskip}{3pt}
\emergencystretch=3em

% ---------- Identification de la fiche -------------------------
\newcommand{\ficheid}{}
\newcommand{\fichetitre}{}
\newcommand{\fiche}[3]{%
  \renewcommand{\ficheid}{#1}\renewcommand{\fichetitre}{#2}%
  \pagestyle{fancy}\fancyhf{}%
  \fancyhead[L]{\small\textsf{Advanced Probabilities -- Fiche #1}}%
  \fancyhead[R]{\small\textsf{#2}}%
  \fancyfoot[C]{\small\thepage}%
  \begin{center}
    {\LARGE\bfseries Fiche #1 -- #2}\par\smallskip
    {\large\itshape (#3)}\par\smallskip
    {\small Cours de V.\,M.~Lambert, M1 -- semestre 1}
  \end{center}\vspace{-2pt}\hrule\vspace{6pt}}

% ---------- Environnements numérotés (un seul compteur) --------
\theoremstyle{definition}
\newtheorem{definition}{Définition}
\renewcommand{\thedefinition}{\ficheid.\arabic{definition}}
\newtheorem{theoreme}[definition]{Théorème}
\newtheorem{proposition}[definition]{Proposition}
\newtheorem{lemme}[definition]{Lemme}
\newtheorem{corollaire}[definition]{Corollaire}
\newtheorem{exemple}[definition]{Exemple}
\theoremstyle{remark}
\newtheorem{remarque}[definition]{Remarque}

% Résultat hors cours (standard, signalé comme tel)
\theoremstyle{definition}
\newtheorem{horscours}[definition]{Hors cours}

% Idée de preuve
\newenvironment{idee}{\par\small\textit{Idée de preuve.}~}{\par}

% ---------- Encadrés -----------------------------------------
\tcbset{fichebox/.style={enhanced,breakable,boxrule=0.6pt,arc=2pt,
  left=5pt,right=5pt,top=3pt,bottom=3pt,fonttitle=\bfseries}}
\newtcolorbox{piege}[1][]{fichebox,colback=red!4,colframe=red!60!black,
  title={Pièges classiques},#1}
\newtcolorbox{liens}[1][]{fichebox,colback=blue!3,colframe=blue!50!black,
  title={Liens entre notions},#1}
\newtcolorbox{methode}[1][Méthode]{fichebox,colback=green!4,colframe=green!40!black,title={#1}}
\newtcolorbox{alerte}[1][]{fichebox,colback=orange!6,colframe=orange!70!black,
  title={Attention à l'énoncé},#1}

% Exercice de TD (fiches de rappels)
\newcounter{exo}
\newenvironment{exo}[1][]{\refstepcounter{exo}\par\medskip
  \noindent{\large\bfseries Exercice~\theexo}\ifx\relax#1\relax\else~{\large\bfseries -- #1}\fi\par\nopagebreak\smallskip}{\par}
\newcommand{\outils}{\textbf{Outils du cours.}~}
\newcommand{\meth}{\textbf{Méthode / astuce.}~}
\newcommand{\indic}{\textbf{Indication.}~}
\newcommand{\rep}{\textbf{Réponse.}~}

% ---------- Notations du cours ---------------------------------
\newcommand{\R}{\mathbb{R}}
\newcommand{\N}{\mathbb{N}}
\newcommand{\Q}{\mathbb{Q}}
\newcommand{\C}{\mathbb{C}}
\newcommand{\PP}{\mathbb{P}}
\newcommand{\QQ}{\mathbb{Q}}
\newcommand{\E}{\mathbb{E}}
\newcommand{\F}{\mathcal{F}}
\newcommand{\G}{\mathcal{G}}
\newcommand{\Hc}{\mathcal{H}}
\newcommand{\Tc}{\mathcal{T}}
\newcommand{\Bc}{\mathcal{B}}
\newcommand{\Cc}{\mathcal{C}}
\newcommand{\Ec}{\mathcal{E}}
\newcommand{\Pc}{\mathcal{P}}
\newcommand{\Nc}{\mathcal{N}}
\newcommand{\LN}{\mathcal{LN}}
\newcommand{\Expo}{\mathcal{E}}
\newcommand{\Poi}{\mathcal{P}}
\newcommand{\Ff}{\mathbb{F}}
\newcommand{\1}{\mathbf{1}}
\newcommand{\dd}{\,\mathrm{d}}
\DeclareMathOperator{\Var}{\mathbb{V}}
\DeclareMathOperator{\Cov}{Cov}
\DeclareMathOperator{\card}{card}
\DeclareMathOperator{\diag}{diag}
\newcommand{\indep}{\perp\!\!\!\perp}
\newcommand{\ps}{\ \PP\text{-p.s.}}
\newcommand{\tendloi}{\xrightarrow[n\to+\infty]{\text{loi}}}
\newcommand{\tendps}{\xrightarrow[n\to+\infty]{\text{p.s.}}}
\newcommand{\tendproba}{\xrightarrow[n\to+\infty]{\PP}}
\newcommand{\tendLp}{\xrightarrow[n\to+\infty]{L^p}}
\newcommand{\tend}{\xrightarrow[n\to+\infty]{}}

%     \fiche{C5}{Vecteurs gaussiens}{Gaussian vectors}
% ==============================================================
\usepackage[T1]{fontenc}
\usepackage[utf8]{inputenc}
\usepackage[french]{babel}
\usepackage{lmodern}
\usepackage[a4paper,margin=1.9cm,top=1.7cm,bottom=1.9cm]{geometry}
\usepackage{amsmath,amssymb,amsthm,mathtools}
\usepackage[inline]{enumitem}
\usepackage{booktabs,array,tabularx}
\usepackage{xcolor}
\usepackage[most]{tcolorbox}
\usepackage{fancyhdr}
\usepackage{hyperref}
\hypersetup{colorlinks=true,linkcolor=blue!50!black,urlcolor=blue!50!black}

\setlist{itemsep=1pt,topsep=3pt,parsep=0pt}
\setlength{\parindent}{0pt}
\setlength{\parskip}{3pt}
\emergencystretch=3em

% ---------- Identification de la fiche -------------------------
\newcommand{\ficheid}{}
\newcommand{\fichetitre}{}
\newcommand{\fiche}[3]{%
  \renewcommand{\ficheid}{#1}\renewcommand{\fichetitre}{#2}%
  \pagestyle{fancy}\fancyhf{}%
  \fancyhead[L]{\small\textsf{Advanced Probabilities -- Fiche #1}}%
  \fancyhead[R]{\small\textsf{#2}}%
  \fancyfoot[C]{\small\thepage}%
  \begin{center}
    {\LARGE\bfseries Fiche #1 -- #2}\par\smallskip
    {\large\itshape (#3)}\par\smallskip
    {\small Cours de V.\,M.~Lambert, M1 -- semestre 1}
  \end{center}\vspace{-2pt}\hrule\vspace{6pt}}

% ---------- Environnements numérotés (un seul compteur) --------
\theoremstyle{definition}
\newtheorem{definition}{Définition}
\renewcommand{\thedefinition}{\ficheid.\arabic{definition}}
\newtheorem{theoreme}[definition]{Théorème}
\newtheorem{proposition}[definition]{Proposition}
\newtheorem{lemme}[definition]{Lemme}
\newtheorem{corollaire}[definition]{Corollaire}
\newtheorem{exemple}[definition]{Exemple}
\theoremstyle{remark}
\newtheorem{remarque}[definition]{Remarque}

% Résultat hors cours (standard, signalé comme tel)
\theoremstyle{definition}
\newtheorem{horscours}[definition]{Hors cours}

% Idée de preuve
\newenvironment{idee}{\par\small\textit{Idée de preuve.}~}{\par}

% ---------- Encadrés -----------------------------------------
\tcbset{fichebox/.style={enhanced,breakable,boxrule=0.6pt,arc=2pt,
  left=5pt,right=5pt,top=3pt,bottom=3pt,fonttitle=\bfseries}}
\newtcolorbox{piege}[1][]{fichebox,colback=red!4,colframe=red!60!black,
  title={Pièges classiques},#1}
\newtcolorbox{liens}[1][]{fichebox,colback=blue!3,colframe=blue!50!black,
  title={Liens entre notions},#1}
\newtcolorbox{methode}[1][Méthode]{fichebox,colback=green!4,colframe=green!40!black,title={#1}}
\newtcolorbox{alerte}[1][]{fichebox,colback=orange!6,colframe=orange!70!black,
  title={Attention à l'énoncé},#1}

% Exercice de TD (fiches de rappels)
\newcounter{exo}
\newenvironment{exo}[1][]{\refstepcounter{exo}\par\medskip
  \noindent{\large\bfseries Exercice~\theexo}\ifx\relax#1\relax\else~{\large\bfseries -- #1}\fi\par\nopagebreak\smallskip}{\par}
\newcommand{\outils}{\textbf{Outils du cours.}~}
\newcommand{\meth}{\textbf{Méthode / astuce.}~}
\newcommand{\indic}{\textbf{Indication.}~}
\newcommand{\rep}{\textbf{Réponse.}~}

% ---------- Notations du cours ---------------------------------
\newcommand{\R}{\mathbb{R}}
\newcommand{\N}{\mathbb{N}}
\newcommand{\Q}{\mathbb{Q}}
\newcommand{\C}{\mathbb{C}}
\newcommand{\PP}{\mathbb{P}}
\newcommand{\QQ}{\mathbb{Q}}
\newcommand{\E}{\mathbb{E}}
\newcommand{\F}{\mathcal{F}}
\newcommand{\G}{\mathcal{G}}
\newcommand{\Hc}{\mathcal{H}}
\newcommand{\Tc}{\mathcal{T}}
\newcommand{\Bc}{\mathcal{B}}
\newcommand{\Cc}{\mathcal{C}}
\newcommand{\Ec}{\mathcal{E}}
\newcommand{\Pc}{\mathcal{P}}
\newcommand{\Nc}{\mathcal{N}}
\newcommand{\LN}{\mathcal{LN}}
\newcommand{\Expo}{\mathcal{E}}
\newcommand{\Poi}{\mathcal{P}}
\newcommand{\Ff}{\mathbb{F}}
\newcommand{\1}{\mathbf{1}}
\newcommand{\dd}{\,\mathrm{d}}
\DeclareMathOperator{\Var}{\mathbb{V}}
\DeclareMathOperator{\Cov}{Cov}
\DeclareMathOperator{\card}{card}
\DeclareMathOperator{\diag}{diag}
\newcommand{\indep}{\perp\!\!\!\perp}
\newcommand{\ps}{\ \PP\text{-p.s.}}
\newcommand{\tendloi}{\xrightarrow[n\to+\infty]{\text{loi}}}
\newcommand{\tendps}{\xrightarrow[n\to+\infty]{\text{p.s.}}}
\newcommand{\tendproba}{\xrightarrow[n\to+\infty]{\PP}}
\newcommand{\tendLp}{\xrightarrow[n\to+\infty]{L^p}}
\newcommand{\tend}{\xrightarrow[n\to+\infty]{}}

%     \fiche{C5}{Vecteurs gaussiens}{Gaussian vectors}
% ==============================================================
\usepackage[T1]{fontenc}
\usepackage[utf8]{inputenc}
\usepackage[french]{babel}
\usepackage{lmodern}
\usepackage[a4paper,margin=1.9cm,top=1.7cm,bottom=1.9cm]{geometry}
\usepackage{amsmath,amssymb,amsthm,mathtools}
\usepackage[inline]{enumitem}
\usepackage{booktabs,array,tabularx}
\usepackage{xcolor}
\usepackage[most]{tcolorbox}
\usepackage{fancyhdr}
\usepackage{hyperref}
\hypersetup{colorlinks=true,linkcolor=blue!50!black,urlcolor=blue!50!black}

\setlist{itemsep=1pt,topsep=3pt,parsep=0pt}
\setlength{\parindent}{0pt}
\setlength{\parskip}{3pt}
\emergencystretch=3em

% ---------- Identification de la fiche -------------------------
\newcommand{\ficheid}{}
\newcommand{\fichetitre}{}
\newcommand{\fiche}[3]{%
  \renewcommand{\ficheid}{#1}\renewcommand{\fichetitre}{#2}%
  \pagestyle{fancy}\fancyhf{}%
  \fancyhead[L]{\small\textsf{Advanced Probabilities -- Fiche #1}}%
  \fancyhead[R]{\small\textsf{#2}}%
  \fancyfoot[C]{\small\thepage}%
  \begin{center}
    {\LARGE\bfseries Fiche #1 -- #2}\par\smallskip
    {\large\itshape (#3)}\par\smallskip
    {\small Cours de V.\,M.~Lambert, M1 -- semestre 1}
  \end{center}\vspace{-2pt}\hrule\vspace{6pt}}

% ---------- Environnements numérotés (un seul compteur) --------
\theoremstyle{definition}
\newtheorem{definition}{Définition}
\renewcommand{\thedefinition}{\ficheid.\arabic{definition}}
\newtheorem{theoreme}[definition]{Théorème}
\newtheorem{proposition}[definition]{Proposition}
\newtheorem{lemme}[definition]{Lemme}
\newtheorem{corollaire}[definition]{Corollaire}
\newtheorem{exemple}[definition]{Exemple}
\theoremstyle{remark}
\newtheorem{remarque}[definition]{Remarque}

% Résultat hors cours (standard, signalé comme tel)
\theoremstyle{definition}
\newtheorem{horscours}[definition]{Hors cours}

% Idée de preuve
\newenvironment{idee}{\par\small\textit{Idée de preuve.}~}{\par}

% ---------- Encadrés -----------------------------------------
\tcbset{fichebox/.style={enhanced,breakable,boxrule=0.6pt,arc=2pt,
  left=5pt,right=5pt,top=3pt,bottom=3pt,fonttitle=\bfseries}}
\newtcolorbox{piege}[1][]{fichebox,colback=red!4,colframe=red!60!black,
  title={Pièges classiques},#1}
\newtcolorbox{liens}[1][]{fichebox,colback=blue!3,colframe=blue!50!black,
  title={Liens entre notions},#1}
\newtcolorbox{methode}[1][Méthode]{fichebox,colback=green!4,colframe=green!40!black,title={#1}}
\newtcolorbox{alerte}[1][]{fichebox,colback=orange!6,colframe=orange!70!black,
  title={Attention à l'énoncé},#1}

% Exercice de TD (fiches de rappels)
\newcounter{exo}
\newenvironment{exo}[1][]{\refstepcounter{exo}\par\medskip
  \noindent{\large\bfseries Exercice~\theexo}\ifx\relax#1\relax\else~{\large\bfseries -- #1}\fi\par\nopagebreak\smallskip}{\par}
\newcommand{\outils}{\textbf{Outils du cours.}~}
\newcommand{\meth}{\textbf{Méthode / astuce.}~}
\newcommand{\indic}{\textbf{Indication.}~}
\newcommand{\rep}{\textbf{Réponse.}~}

% ---------- Notations du cours ---------------------------------
\newcommand{\R}{\mathbb{R}}
\newcommand{\N}{\mathbb{N}}
\newcommand{\Q}{\mathbb{Q}}
\newcommand{\C}{\mathbb{C}}
\newcommand{\PP}{\mathbb{P}}
\newcommand{\QQ}{\mathbb{Q}}
\newcommand{\E}{\mathbb{E}}
\newcommand{\F}{\mathcal{F}}
\newcommand{\G}{\mathcal{G}}
\newcommand{\Hc}{\mathcal{H}}
\newcommand{\Tc}{\mathcal{T}}
\newcommand{\Bc}{\mathcal{B}}
\newcommand{\Cc}{\mathcal{C}}
\newcommand{\Ec}{\mathcal{E}}
\newcommand{\Pc}{\mathcal{P}}
\newcommand{\Nc}{\mathcal{N}}
\newcommand{\LN}{\mathcal{LN}}
\newcommand{\Expo}{\mathcal{E}}
\newcommand{\Poi}{\mathcal{P}}
\newcommand{\Ff}{\mathbb{F}}
\newcommand{\1}{\mathbf{1}}
\newcommand{\dd}{\,\mathrm{d}}
\DeclareMathOperator{\Var}{\mathbb{V}}
\DeclareMathOperator{\Cov}{Cov}
\DeclareMathOperator{\card}{card}
\DeclareMathOperator{\diag}{diag}
\newcommand{\indep}{\perp\!\!\!\perp}
\newcommand{\ps}{\ \PP\text{-p.s.}}
\newcommand{\tendloi}{\xrightarrow[n\to+\infty]{\text{loi}}}
\newcommand{\tendps}{\xrightarrow[n\to+\infty]{\text{p.s.}}}
\newcommand{\tendproba}{\xrightarrow[n\to+\infty]{\PP}}
\newcommand{\tendLp}{\xrightarrow[n\to+\infty]{L^p}}
\newcommand{\tend}{\xrightarrow[n\to+\infty]{}}

\begin{document}
\fiche{TD5}{Processus, martingales, temps d'arrêt, Poisson}{Random Process, Martingale, Stopping Time and Poisson Processes}

\textit{À lire avant de chercher le TD. Renvois : « C8.6 » = fiche C8, énoncé n°6. Fiches associées : \textbf{C8}, \textbf{C6}, \textbf{C9}, \textbf{C3}.}

\begin{exo}[Temps d'arrêt]
\outils Définition d'un temps d'arrêt (C8.6) ; temps d'entrée (C8.7 et sa remarque) ; tribu (C1.1).

\meth Écrire $\{\tau\le t\}$ à l'aide d'événements déjà dans $\F_t$.

\indic 1. $\{t_0\le t\}\in\{\emptyset,\Omega\}$. 2. $\{\tau_1\vee\tau_2\le t\}=\{\tau_1\le t\}\cap\{\tau_2\le t\}$ et
$\{\tau_1\wedge\tau_2\le t\}=\{\tau_1\le t\}\cup\{\tau_2\le t\}$.
3. Avec $d(x,F)=\inf_{y\in F}|x-y|$ (continue, nulle exactement sur $F$ fermé) et les trajectoires continues :
$\{\tau\le t\}=\bigl\{\inf_{s\in([0,t]\cap\Q)\cup\{t\}}d(X_s,F)=0\bigr\}$. C'est un infimum \emph{dénombrable} de v.a.\ $\F_t$-mesurables.
\end{exo}

\begin{exo}[Martingales multiplicatives et logarithmiques]
\outils Martingale en temps discret (C8.2, C8.4) ; indépendance $X_{n+1}\indep\F_n$ (C3.11) ; $\E(\cdot\mid\G)=\E(\cdot)$ si indépendance
(C6.5, point 7) ; sortir ce qui est connu (C6.5, point 4).

\meth $\ln P_n=\sum_{k\le n}\ln(1+X_k)$ est une somme de v.a.\ i.i.d.\ intégrables : on la recentre. $P_{n+1}=P_n(1+X_{n+1})$ avec
$P_n$ $\F_n$-mesurable. Vérifier aussi l'intégrabilité ($P_n$ est un produit de v.a.\ indépendantes intégrables, C3.15).

\rep 1. $a=\E\bigl(\ln(1+X_1)\bigr)$. 2. $b=\E(1+X_1)=1+\E X_1$ (supposé non nul).
\end{exo}

\begin{exo}[Casino : on ne bat pas un jeu équitable]
\begin{alerte}
L'hypothèse « $\{\tau=i\}$ indépendant de $\F_i$ » est dégénérée si $\tau$ est un temps d'arrêt : comme $\{\tau=i\}\in\F_i$, cet événement
serait indépendant de lui-même, donc $\PP(\tau=i)=\PP(\tau=i)^2\in\{0,1\}$ et $\tau$ serait déterministe.
% Vérifié : {tau=i} dans F_i et independant de F_i => P(tau=i) dans {0,1}.
Deux lectures donnent le résultat sans dégénérescence : (a) $\tau$ temps d'arrêt borné par $N$ : écrire
$S_\tau=\sum_{i\le N}X_i\1_{\{\tau\ge i\}}$ avec $\{\tau\ge i\}=\{\tau\le i-1\}^c\in\F_{i-1}$, indépendant de $X_i$ ;
(b) $\tau$ indépendant de toute la suite $(X_i)$ (arrêt décidé « de l'extérieur »).
\end{alerte}
\outils Martingale (C8.2) ; sortir ce qui est connu et indépendance (C6.5) ; temps d'arrêt (C8.6) ; théorème d'arrêt (C8.10).

\meth
\begin{itemize}
  \item $\F_n$ = information après $n$ parties. $\E(S_{n+1}\mid\F_n)=S_n+\E(X_{n+1})$.
  \item Stratégie : la mise $H_n$ doit être décidée \textbf{avant} la $n$-ième partie, donc $\F_{n-1}$-mesurable (\emph{prévisible}),
  et telle que $H_nX_n\in L^1$ (par ex.\ $H_n$ bornée). Alors on peut sortir $H_{n+1}$ de $\E(H_{n+1}X_{n+1}\mid\F_n)$.
  \item Arrêt : $\{\tau=n\}\in\F_n$ (décider d'arrêter avec l'information disponible). Calculer
  $\E(S_\tau)=\sum_i\E(S_i\1_{\{\tau=i\}})$ ($\tau$ borné : somme finie) et utiliser l'hypothèse de l'énoncé ; ou, plus solidement, la décomposition (a) de l'encadré.
\end{itemize}

\rep $\E(S_{n+1}\mid\F_n)=S_n$ ; $\E(\tilde S_{n+1}\mid\F_n)=\tilde S_n$ (toute stratégie prévisible reste équitable) ; $\E(S_\tau)=0$
(cohérent avec C8.10).
\end{exo}

\begin{exo}[Jeu du capital partagé]
\outils Martingale bornée et convergence p.s.\ (C8.11, hors cours) ; convergence dominée ; espérance totale (C6.5) ; $\liminf$ (C3.5) ;
Borel--Cantelli 2 (C3.6).

\meth
\begin{enumerate}
  \item Appliquer la relation de l'énoncé avec $f=\mathrm{id}$ ; $X_n\in[0,1]$ donne l'intégrabilité.
  \item Appliquer avec $f(x)=x^2$, prendre l'espérance ; passer à la limite par convergence dominée ($0\le X_n\le1$) dans
  $\E(X_{n+1}^2)=\frac14\E(3X_n^2+X_n)$ et dans $\E(X_n)=p$.
  \item $W(1-W)\ge0$ d'espérance nulle, donc nul p.s.
  \item Gain : $X_{n+1}=X_n+\frac{1-X_n}2$, donc $Y_n=1$ ; perte : $X_{n+1}=\frac{X_n}2$, donc $Y_n=0$. D'où $Y_n=\1_{\{\text{gain à }n+1\}}$.
  \item $Y_n=2X_{n+1}-X_n\to2Z-Z=Z$ p.s. ; une suite à valeurs dans $\{0,1\}$ qui converge est stationnaire.
  Pour l'indépendance : si les $Y_n$ étaient indépendantes, que dirait BC2 appliqué aux $H_n$ ?
\end{enumerate}

\rep 2. $\E(Z^2)=\E(Z)=p$. 3. $Z\sim\mathcal B(p)$. 4. $\E(\1_{\{Y_n=1\}}\mid\F_n)=X_n$, $\E(\1_{\{Y_n=0\}}\mid\F_n)=1-X_n$, donc $Y_n\sim\mathcal B(p)$.
5. $\liminf G_n=\{Z=1\}$, de probabilité $p$, et $\liminf H_n=\{Z=0\}$, de probabilité $1-p$ ; les $Y_n$ ne sont \emph{pas} indépendantes (sinon BC2 donnerait $\PP(\liminf G_n)=0$).
6. Avec probabilité $p$, le joueur finit par gagner toutes les parties et capte tout le capital ; avec probabilité $1-p$, il perd tout.
\end{exo}

\begin{exo}[Carré d'une sous-martingale positive]
\outils Jensen conditionnel (C6.5, point 8) ; monotonie (C6.5, point 3) ; sous-martingale (C8.2).

\meth $\E(M_t^2\mid\F_s)\ge\bigl(\E(M_t\mid\F_s)\bigr)^2\ge M_s^2$. La seconde inégalité utilise $\E(M_t\mid\F_s)\ge M_s\ge0$ et la
croissance de $x\mapsto x^2$ sur $\R_+$. Intégrabilité : $M_t\in L^2$.
\end{exo}

\begin{exo}[Fonction convexe croissante d'une sous-martingale]
\outils Mêmes outils qu'à l'exercice 5.

\meth $\E(\Phi(M_t)\mid\F_s)\ge\Phi(\E(M_t\mid\F_s))\ge\Phi(M_s)$ : Jensen, puis croissance de $\Phi$.
Lien : l'exercice 5 est le cas $\Phi(x)=(x^+)^2$ (convexe croissante sur $\R$, égale à $x^2$ sur $\R_+$, où vit $M$).
\end{exo}

\begin{exo}[Martingale fermée]
\outils Tour (C6.5, point 6) ; Jensen avec $|\cdot|$ ; martingale (C8.2, C8.3).

\meth Pour $s\le t$, $\F_s\subset\F_t$ : $\E(X_t\mid\F_s)=\E(\E(Y\mid\F_t)\mid\F_s)=\E(Y\mid\F_s)=X_s$. Intégrabilité :
$\E|X_t|\le\E\bigl(\E(|Y|\mid\F_t)\bigr)=\E|Y|$. Adaptation : par définition de l'espérance conditionnelle.
\end{exo}

\begin{exo}[Poisson compensé]
\outils Déf.\ du processus de Poisson (C9.4) ; indépendance et $\E(\cdot\mid\G)$ (C6.5, point 7) ; regroupement par paquets (C3.11).

\meth $\F_n=\sigma(N_0,\dots,N_n)=\sigma(N_1-N_0,\dots,N_n-N_{n-1})$ ; l'accroissement $N_{n+1}-N_n$ est indépendant de ces
accroissements, et a même loi que $N_1\sim\Poi(\lambda)$. Donc $\E(N_{n+1}-N_n\mid\F_n)=\lambda$. Ne pas oublier adaptation et intégrabilité.
\end{exo}

\begin{exo}[Attente résiduelle]
\outils Thm C9.6 (c'est exactement l'énoncé du cours) ; dictionnaire C9.3.

\meth $\{W_1>t\}=\{N_{s+t}-N_s=0\}$, puis stationnarité des accroissements.

\rep $W_1\sim\Expo(\lambda)$ ; $T_1\sim\Expo(\lambda)$ (cas $s=0$).
\end{exo}

\begin{exo}[Instants de saut et statistiques d'ordre]
\outils Densité produit des $X_i$ indépendantes (C3.16) ; changement de variables linéaire (C2.10, version multidimensionnelle) ; dictionnaire C9.3 ;
densités conditionnelles (C6.9).

\meth 1. $(T_1,\dots,T_4)=A(X_1,\dots,X_4)$ avec $A$ triangulaire inférieure de 1, $|\det A|=1$ ; image du domaine $\R_+^{*4}$ :
$\{0<t_1<t_2<t_3<t_4\}$. 2. $\{N_t=3\}=\{T_3\le t<T_4\}$ : intégrer la densité de la question 1 en $t_4$ sur $]t,+\infty[$, puis diviser
par $\PP(N_t=3)$.

\rep 1. $f(t_1,\dots,t_4)=\lambda^4e^{-\lambda t_4}\1_{0<t_1<t_2<t_3<t_4}$.
2. $\dfrac{3!}{t^3}\1_{0<t_1<t_2<t_3<t}$ : loi des statistiques d'ordre de 3 uniformes indépendantes sur $[0,t]$.
\end{exo}

\begin{exo}[Le paradoxe de l'autobus]
\outils Dictionnaire C9.3 ; Thm C9.6 ; accroissements indépendants et stationnaires (C9.4) ; critère des fonctions de répartition pour
l'indépendance (C3.9) ; convergence en loi par fonctions de répartition (C4.7).

\meth Pour $0\le z\le t$ et $w\ge0$, l'événement $\{Z_t\ge z,\ W_t>w\}$ signifie qu'il n'y a aucune arrivée dans $]t-z,t+w]$
(avec $Z_t=t$ s'il n'y a eu aucune arrivée avant $t$). Sa probabilité factorise. Alternative suggérée par l'énoncé : décomposer selon
$\{N_t=n\}$ et utiliser la densité d'Erlang de $T_n$.

\rep 1. $\PP(Z_t\ge z,W_t>w)=e^{-\lambda z}e^{-\lambda w}$ pour $0\le z\le t$, $w\ge0$ (et $0$ si $z>t$) ; d'où la fonction de répartition $F(z,w)=(1-e^{-\lambda z})(1-e^{-\lambda w})$ pour $0\le z<t$, $w\ge0$, et $1-e^{-\lambda w}$ pour $z\ge t$ : $Z_t\indep W_t$ et $W_t\sim\Expo(\lambda)$.
2. $Z_t$ a la loi de $\min(T_1,t)$ (atome $e^{-\lambda t}$ en $t$) ; sa fonction de répartition tend vers celle de $\Expo(\lambda)$.
3. $\E(Z_t+W_t)=\frac{1-e^{-\lambda t}}{\lambda}+\frac1\lambda\xrightarrow[t\to\infty]{}\frac2\lambda$ : l'intervalle qui contient $t$ est en moyenne
\emph{deux fois} plus long qu'un intervalle typique ($\frac1\lambda$), car on tombe plus souvent dans les grands intervalles.
4. Attente moyenne $\E(W_t)=\frac1\lambda$, et non $\frac1{2\lambda}$.
\end{exo}

\section*{Synthèse : type de question $\to$ méthode}
\renewcommand{\arraystretch}{1.3}
\begin{tabularx}{\linewidth}{@{}>{\raggedright\arraybackslash}p{0.36\linewidth}X@{}}
\toprule
Type de question & Méthode\\\midrule
Montrer que $\tau$ est un temps d'arrêt & Écrire $\{\tau\le t\}$ (ou $\{\tau=n\}$) avec des événements de $\F_t$ ; opérations dénombrables (C8.6).\\
Montrer que $(X_n)$ est une martingale & Adapté + intégrable + $\E(X_{n+1}\mid\F_n)=X_n$ : sortir le connu, l'indépendant devient une espérance (C8.4, C6.5).\\
Trouver la constante qui rend une suite martingale & Calculer $\E(X_{n+1}\mid\F_n)$ en fonction du paramètre, puis identifier.\\
Stratégies, mises & Mise prévisible et bornée $\Rightarrow$ la transformée de martingale est une martingale.\\
Sous-martingale $\Phi(M)$ & Jensen conditionnel + monotonie de $\Phi$ (C6.5).\\
Limite d'une martingale bornée & Convergence p.s.\ (C8.11) + convergence dominée pour passer les espérances à la limite.\\
Indépendance d'une suite ? & Confronter à Borel--Cantelli 2 (C3.6).\\
Poisson : lois d'attente & Traduire en $\{N_{s+t}-N_s=0\}$ ; stationnarité et indépendance des accroissements (C9.3, C9.6).\\
Densité des instants de saut & Changement de variables linéaire $T=AX$ ; conditionner par $\{N_t=n\}=\{T_n\le t<T_{n+1}\}$.\\\bottomrule
\end{tabularx}

\end{document}
